\chapter{Einführung}
\label{ch:Einfuehrung}


\section{Das \glsentryshort{frost}-Labor}
\label{sec:frostlabor}

Das \gls{frost}-Labor ist ein studentisches Team der \gls{frauas}, welches Robotersysteme entwickelt. Gezielt werden Projekte für Bachelor- und Masterarbeiten ausgeschrieben um an der \gls{erc} teilzunehmen, ein internationaler Wettbewerb, bei dem Teams Marsrover bauen und in simulierten Missionen antreten. Aufgeteilt in sieben verschiedene Aufgaben ein selbstgebauter Marsrover, begleitet von einer Drohne, treten 25 Teams gegeneinander an.

\subsection{Die \glsentryshort{daedalus}-Drohne}
\label{sec:daedalusdrone}

Die Drohne \gls{daedalus} wurde aus Komponenten von einem Holybro X500 V2 kit gebaut. Bei der Drohennaufgabe des \gls{erc}-Wettbewerbs wird aus der Luft eine Fläche abgeflogen und mit Computer-Vision nach Proben gescannt, welche für eine spätere Untersuchung von seitens des Rovers markiert werden. Dabei ist ein stabiler und ruhiger Flug wichtig, bei der die Ausrichtung der Motoren, das Gewicht, Kameraposition und Vibrationsvermeidung eine große Rolle spielen.




% ...



%Am Ende des Einführungskapitels ein kurzer Überblick über die kommenden Kapitel in einem Absatz (Beschreibung der Struktur der Thesis):

%...